\documentclass[11pt]{article}
\usepackage[utf8]{inputenc}
\usepackage[T1]{fontenc}
\usepackage[margin=1in]{geometry}
\usepackage{hyperref}
\usepackage{enumitem}
\hypersetup{colorlinks=true,linkcolor=black,urlcolor=blue,citecolor=blue}
\title{Intro Typing and Simple Circuits}
\author{Piter Garcia}
\date{March 20, 2026}
\begin{document}
\maketitle
\section*{Packet Use}
This printable lesson packet is generated from the paired Markdown guide. The Markdown version remains the clickable, neurodivergent-friendly working copy; this PDF carries the same lesson substance in a formal handout format. Evidence claims in this packet are based on transcript-derived lesson notes and the cleaned transcript files in this workspace.

\section*{Quick Scan}
\begin{itemize}[leftmargin=*]
\item \textbf{Grade:} Kindergarten
\item \textbf{Grade Hub:} Notes [../kindergarten-Notes.md]
\item \textbf{Schedule Reference:} Spring2026\_TeachingSchedule\_Derek.md [../../school/Spring2026\_TeachingSchedule\_Derek.md]
\item \textbf{Workbook Sheet:} \texttt{K} in \texttt{Teaching\_Placement-IGNITE Curriculum\_ GCSD25-26.xlsx}
\item \textbf{Lesson Focus:} Short keyboarding routines paired with a teacher-guided introduction to simple circuits using Play-Doh and classroom-safe materials.
\item \textbf{CS Alignment Strength:} Foundational CS + physical computing bridge
\item \textbf{LaTeX Source:} intro-typing-and-simple-circuits.tex [./intro-typing-and-simple-circuits.tex]
\item \textbf{Bib File:} intro-typing-and-simple-circuits.bib [./intro-typing-and-simple-circuits.bib]
\item \textbf{Artifacts Folder:} artifacts/ [./artifacts/]
\end{itemize}

\section*{Lesson Identity}
\begin{itemize}[leftmargin=*]
\item \textbf{Likely Class Blocks:} Day 1 -- 12:25-1:15 -- Justinger, Day 2 -- 12:25-1:15 -- Kesys, Day 4 -- 12:25-1:15 -- Pum
\item \textbf{Main Lesson Family:} Intro Typing and Simple Circuits
\item \textbf{Primary Evidence Base:} Transcript-derived notes plus source transcripts from this teaching-placement workspace.
\item \textbf{Primary External Source Type:} Official curriculum/provider materials.
\end{itemize}

\section*{What We Are Teaching}
\begin{itemize}[leftmargin=*]
\item home-row style finger placement and space-bar recognition
\item following a step sequence before receiving the hands-on science task
\item identifying battery, bulb, and connection points in a circuit activity
\item understanding that a working system depends on complete connections
\end{itemize}

\section*{CS Standards Alignment}
\begin{itemize}[leftmargin=*]
\item \textbf{1A-CS-02}: Use appropriate terminology to describe common physical components of computing systems. Why it fits here: Students name and use keyboard parts and circuit components such as battery and bulb.
\item \textbf{1A-AP-08}: Model daily processes by creating and following algorithms. Why it fits here: Both the typing warm-up and the circuit build depend on following steps in order.
\item \textbf{1A-AP-11}: Decompose the steps needed to solve a problem into a precise sequence of instructions. Why it fits here: Teacher instruction explicitly breaks the task into one step at a time before students try independently.
\end{itemize}

\section*{NYS / Local Curriculum Alignment}
\begin{itemize}[leftmargin=*]
\item Aligned to the local IGNITE kindergarten workbook sequence through keyboard entry routines, safe material handling, and early engineering-style problem solving.
\item Supports New York-facing expectations for digital fluency and hands-on scientific observation by connecting device use with simple cause-and-effect investigation.
\item Bridges computing systems with early physical computing by showing that both keyboard routines and circuits depend on correct inputs and complete connections.
\end{itemize}

\section*{Lesson Content and Concepts}
\begin{itemize}[leftmargin=*]
\item sequence
\item input
\item system parts
\item closed circuit
\item debugging by checking connections
\end{itemize}

\section*{Evidence / Proof of Instruction}
\begin{itemize}[leftmargin=*]
\item \textbf{Transcript-derived note:} 260305 Kindergarten Justinger Typing Circuits [../../transcript-derived-lessons/260305-kindergarten-justinger-typing-circuits.md] The lesson note shows a two-part structure: typing first, then Play-Doh circuits after students complete the warm-up.
\item \textbf{Transcript-derived note:} 260305 Kindergarten Justinger Circuits Summary [../../transcript-derived-lessons/260305-kindergarten-justinger-circuits-summary.md] The planning summary confirms that the circuit activity was intentionally paired with the typing session, not improvised.
\item \textbf{Source transcript:} 260305 Kindergarten Justinger Typing Circuits [../../transcripts/260305-kindergarten-justinger-typing-circuits.md] The transcript records direct modeling of F, J, and space, followed by explicit directions to split Play-Doh, avoid mixing colors, and build with provided materials.
\item \textbf{Likely student proof:} Students completing the typing prompt, assembling a working or near-working circuit, and verbally naming parts count as direct evidence of instruction.
\end{itemize}

\section*{Assessment Evidence}
\begin{itemize}[leftmargin=*]
\item Student finds required keys with prompting reduced over time.
\item Student can identify at least one circuit part and explain that contact matters.
\item Student participates safely in the material routine and attempts the build sequence.
\end{itemize}

\section*{UDL and Inclusion Supports}
\subsection*{Multiple Representation}
\begin{itemize}[leftmargin=*]
\item Show keys visually and physically point to them on the keyboard.
\item Model circuit parts one at a time before handing out all materials.
\item Use color and spatial language consistently: here, there, together, connected.
\end{itemize}

\subsection*{Multiple Action / Expression}
\begin{itemize}[leftmargin=*]
\item Allow partner-supported keyboarding and hand-over-hand support when needed.
\item Offer pre-portioned materials and physically larger examples.
\item Accept explanation through pointing, showing, or brief oral response.
\end{itemize}

\subsection*{Multiple Engagement}
\begin{itemize}[leftmargin=*]
\item Use the science activity as a motivating payoff for the typing warm-up.
\item Keep the typing segment short and success-based.
\item Make the transition clear so students know when one task ends and the next begins.
\end{itemize}

\section*{How We Are Already Making It Inclusive}
\begin{itemize}[leftmargin=*]
\item The transcript shows strong teacher modeling, immediate feedback, and a high-support release into practice.
\item Students were given a clear transition reward structure, which helped sustain attention through the warm-up.
\item The hands-on circuit task created a non-text entry point for students who struggle with abstract explanation.
\end{itemize}

\section*{How to Improve Inclusion Next Time}
\begin{itemize}[leftmargin=*]
\item Add visual step mats for both the keyboard task and the circuit build.
\item Prepare one complete exemplar circuit at each table for students who need a concrete anchor.
\item Add sentence stems such as “It worked because…” and “We need to fix…”
\item Create a quieter typing option or reduced-key version for students who fatigue quickly.
\end{itemize}

\section*{Materials and Setup}
\begin{itemize}[leftmargin=*]
\item Student devices
\item Typing activity or typing game
\item Play-Doh
\item Battery pack or coin battery demo
\item Bulb/LED and connectors
\item Projected model
\end{itemize}

\section*{Teacher Moves / Script / Routines}
\begin{itemize}[leftmargin=*]
\item Model the typing task with only the needed keys visible in student attention.
\item Pause before the circuit release to restate one safety rule and one build rule.
\item Circulate first to tables that need connection help rather than re-explaining to the whole group.
\item Close by asking students to show the part that made the system work.
\end{itemize}

\section*{Common Errors and Debugging}
\begin{itemize}[leftmargin=*]
\item Students may confuse keyboard keys or lose left/right hand positioning.
\item Students may mix Play-Doh colors before building and lose structural clarity.
\item Circuit failures will often come from incomplete contact or reversed/loose placement rather than misunderstanding the concept.
\end{itemize}

\section*{Extensions and Simplifications}
\begin{itemize}[leftmargin=*]
\item Extension: have students compare two circuit builds and say which connection is stronger.
\item Simplification: use a teacher demo circuit if student independence is too high a load.
\item Extension: connect typing “input” to the idea that circuits also need the right path to work.
\end{itemize}

\section*{Related Local Materials}
\begin{itemize}[leftmargin=*]
\item No direct local lesson packet linked yet.
\end{itemize}

\section*{Artifacts Checklist}
\begin{itemize}[leftmargin=*]
\item Compiled lesson PDF in \texttt{artifacts/}
\item At least one screenshot, student example, or setup photo
\item One short assessment or observation record
\item Updated standards/source bibliography once the \texttt{.bib} file is revised
\item Any teacher reflection or revision notes after reteaching
\end{itemize}

\section*{Selected Official Sources}
Selected official sources that informed this packet are listed below.
ocite{*}

\bibliographystyle{plain}
\bibliography{intro-typing-and-simple-circuits}
\end{document}
